\documentclass[11pt,a4paper]{article}

\usepackage[utf8]{inputenc}
\usepackage[T1]{fontenc}
\usepackage[margin=3cm,top=3cm,bottom=3cm]{geometry}
\usepackage{listings}
\usepackage{booktabs}
\usepackage{amsmath}
\usepackage{amssymb}
\usepackage{abstract}
\usepackage{fancyhdr}
\usepackage{float}
\usepackage{titling}
\usepackage{authblk}
\usepackage{siunitx}

\lstset{
	basicstyle=\ttfamily\footnotesize,
	breaklines=true,
	breakatwhitespace=false,
	tabsize=8,
	showstringspaces=false,
	numbers=left,
	numberstyle=\tiny,
	numbersep=5pt,
	xleftmargin=10pt,
	frame=lines,
	framerule=0.4pt,
	captionpos=b,
	aboveskip=6pt,
	belowskip=6pt,
}

\pagestyle{fancy}
\fancyhf{}
\lhead{\small\textit{Dispatch Overhead in POSIX Shell: An Empirical Study}}
\rhead{\small\textit{Gaydardzhiev, 2026}}
\cfoot{\thepage}
\renewcommand{\headrulewidth}{0.4pt}
\renewcommand{\footrulewidth}{0pt}

\renewcommand{\abstractnamefont}{\normalfont\small\bfseries}
\renewcommand{\abstracttextfont}{\normalfont\small}
\setlength{\absleftindent}{0pt}
\setlength{\absrightindent}{0pt}

\pretitle{%
	\noindent\rule{\linewidth}{1.2pt}\\[4pt]
	\begin{center}\Large\bfseries
}
\posttitle{\end{center}}
\preauthor{\begin{center}\normalsize}
\postauthor{\end{center}}
\predate{\begin{center}\small}
\postdate{%
	\end{center}\vspace{2pt}
	\noindent\rule{\linewidth}{0.6pt}
}

\begin{document}

\title{Dispatch Overhead in POSIX Shell:\\[4pt]
\large An Empirical Measurement of \texttt{if}/\texttt{[}
versus \texttt{case} Across Shell Implementations}

\author{Ivan Gaydardzhiev}
\affil{\small\textit{Independent Research} \qquad 2026}
\date{}

\maketitle

\begin{abstract}
\noindent
A companion paper argued on structural grounds that \texttt{if~[~]}
is a category error for value-based dispatch in POSIX shell: it routes
a value through a command evaluator, incurring an encoding round-trip
that \texttt{case} does not require.
This paper measures the cost of that round-trip empirically.
Three shell configurations were benchmarked on an ARM Cortex-A53
octa-core system running 32-bit Arch Linux: GNU Bash~5.2.37,
a self-compiled \textsc{dash} (armv8l), and the Pacman-distributed
\textsc{dash} binary.
Each configuration executed a controlled dispatch loop of
100\,000 iterations under both \texttt{if}/\texttt{[} and
\texttt{case}, with arithmetic work held constant.
The \texttt{case} construct outperformed \texttt{if}/\texttt{[} in
every run across all three shells, with mean speedup ratios of
$2.33\times$, $2.35\times$, and $3.09\times$ respectively.
Standard deviations were low and ratios consistent, indicating that
the overhead is structural and not a measurement artefact.
The most significant finding is directional: the faster the shell
implementation, the larger the relative overhead of \texttt{if}/\texttt{[}.
Optimising the shell does not close the gap; it widens it.
\end{abstract}

\vspace{6pt}
\noindent\rule{\linewidth}{0.4pt}
\vspace{4pt}

\section{Introduction}

The structural argument for preferring \texttt{case} over
\texttt{if}/\texttt{[} for value-based dispatch in POSIX shell
has been made elsewhere~\cite{ecnb2026}.
That argument does not require empirical support: the round-trip
is mandated by the specification, and its cost is a property of
the idiom, not of any particular implementation.
Any conformant shell must execute the guard of an \texttt{if}
statement as a command and inspect its exit status.
That requirement cannot be optimised away.

This paper nevertheless measures the cost, for two reasons.

First, a reader who accepts the structural argument on its own terms
may still want to know the order of magnitude.
The difference between a theoretically unnecessary step and a
practically significant one is not always obvious from first
principles.

Second, the measurement reveals a non-obvious property of the
overhead: it does not diminish as the shell implementation improves.
The ratio between \texttt{if}/\texttt{[} and \texttt{case} is
stable or larger in the faster shell.
This is consistent with the structural argument: the round-trip
cost is irreducible, so a faster shell reduces both times
proportionally, leaving the ratio intact or worsened. But it is
not a result one would confidently predict without measuring.

\section{Methodology}

\subsection{Hardware and System}

All measurements were taken on a single machine: a MediaTek MT6765
(Helio G85) system-on-chip with an octa-core ARM Cortex-A53
processor, revision~4, running at nominally 26.08~BogoMIPS per core.
The system ran 32-bit Arch Linux ARM in armv8l compatibility mode,
kernel reporting \texttt{ARMv8 Processor rev~4 (v8l)}.
Physical memory was \SI{2869}{\mega\byte}.
At the time of measurement, approximately \SI{1048}{\mega\byte} was
available; the system was under normal interactive load.

The ARM Cortex-A53 is an in-order pipeline with no out-of-order
execution.
This makes timing more deterministic than a superscalar desktop core
and reduces the variance that would require a larger sample to
characterise.

\subsection{Shell Configurations}

Three shell configurations were benchmarked:

\begin{enumerate}
  \item \textbf{Bash 5.2.37}, GNU Bash release
        \texttt{(armv7l-unknown-linux-gnueabihf)}, invoked as
        \texttt{bash compare.sh}.

  \item \textbf{dash-armv8l}, a \textsc{dash} binary compiled by the
        author from source for the armv8l target, invoked as
        \texttt{./dash-armv8l compare.sh}.

  \item \textbf{dash (Pacman)}, the \textsc{dash} binary distributed
        by the Arch Linux ARM package manager, invoked as
        \texttt{/usr/bin/dash compare.sh}.
\end{enumerate}

The two \textsc{dash} configurations represent different compilation
paths of the same interpreter.
Their differing absolute times reflect compiler and optimisation
differences; their differing ratios are the analytically interesting
result.

\subsection{Benchmark Design}

The benchmark script, \texttt{compare.sh}, defines two functions:
\texttt{ftestif} and \texttt{ftestcase}.
Both iterate a loop of 100\,000 steps.
In each iteration, both functions compute:

\begin{lstlisting}[numbers=none,frame=none,aboveskip=2pt,belowskip=2pt]
r=$((i % 6))
\end{lstlisting}

\noindent
exactly once, assigning the result to \texttt{r}.
The arithmetic work is therefore identical between the two functions.
The only variable is how \texttt{r} is used for dispatch.

\texttt{ftestif} dispatches via \texttt{if}/\texttt{[}:

\begin{lstlisting}[numbers=none,frame=none,aboveskip=2pt,belowskip=2pt]
if   [ "$r" -eq 0 ]; then
    result="even and divisible by three"
elif [ "$r" -eq 2 ] || [ "$r" -eq 4 ]; then
    result="even"
elif [ "$r" -eq 3 ]; then
    result="divisible by three"
else
    result="odd"
fi
\end{lstlisting}

\texttt{ftestcase} dispatches via \texttt{case}:

\begin{lstlisting}[numbers=none,frame=none,aboveskip=2pt,belowskip=2pt]
case $r in
    0)   result="even and divisible by three" ;;
    2|4) result="even"                        ;;
    3)   result="divisible by three"          ;;
    *)   result="odd"                         ;;
esac
\end{lstlisting}

Timing used \texttt{date~+\%s\%N} before and after each loop.
Each configuration was run four times in sequence without rebooting,
producing four paired measurements of \texttt{if} time and
\texttt{case} time per shell.
No CPU pinning or process priority adjustment was applied; the
measurements reflect conditions representative of normal use.

\subsection{A Note on \texttt{if}/\texttt{[} Invocation Count}

In the worst case, a value that does not match the first or second
guard, \texttt{ftestif} invokes \texttt{[} up to four times per
iteration: once for the \texttt{if} guard, once for the \texttt{[}
inside the \texttt{||} disjunction of the first \texttt{elif},
and once each for the remaining \texttt{elif} guards.
\texttt{ftestcase} invokes no external or built-in command for
dispatch at any point.
The benchmark does not attempt to normalise per-invocation cost;
it measures the total cost of expressing the dispatch problem in each
idiom as a programmer would naturally write it.

\section{Results}

Raw timing data across all four runs per shell are presented in
Table~\ref{tab:raw}.
Summary statistics (mean, standard deviation, ratio, and mean
absolute speedup) are presented in Table~\ref{tab:summary}.

\begin{table}[H]
\centering
\caption{Raw timing data (nanoseconds), four runs per shell}
\label{tab:raw}
\small
\begin{tabular}{@{}llrr@{}}
\toprule
Shell & Run & \texttt{if} time (ns) & \texttt{case} time (ns) \\
\midrule
dash-armv8l  & 1 & 7\,007\,660\,624 & 3\,134\,303\,622 \\
(self-compiled) & 2 & 7\,098\,621\,941 & 3\,011\,229\,225 \\
             & 3 & 7\,097\,430\,018 & 3\,009\,546\,762 \\
             & 4 & 7\,043\,410\,166 & 2\,969\,579\,374 \\
\midrule
bash 5.2.37  & 1 & 6\,963\,661\,004 & 2\,953\,720\,680 \\
(armv7l)     & 2 & 7\,047\,592\,090 & 3\,012\,825\,302 \\
             & 3 & 6\,934\,439\,233 & 2\,974\,111\,913 \\
             & 4 & 7\,058\,227\,783 & 2\,986\,966\,068 \\
\midrule
dash         & 1 & 2\,506\,996\,943 & 818\,304\,389 \\
(Pacman)     & 2 & 2\,518\,489\,405 & 800\,539\,080 \\
             & 3 & 2\,497\,728\,250 & 808\,197\,311 \\
             & 4 & 2\,493\,949\,788 & 811\,478\,312 \\
\bottomrule
\end{tabular}
\end{table}

\begin{table}[H]
\centering
\caption{Summary statistics across four runs per shell}
\label{tab:summary}
\small
\begin{tabular}{@{}lrrrrrr@{}}
\toprule
Shell &
  \multicolumn{2}{c}{\texttt{if} (s)} &
  \multicolumn{2}{c}{\texttt{case} (s)} &
  Ratio &
  Speedup (s) \\
\cmidrule(lr){2-3}\cmidrule(lr){4-5}
& Mean & SD (ms) & Mean & SD (ms) & & \\
\midrule
dash-armv8l   & 7.0618 & 44.33 & 3.0312 & 71.40 & 2.330 & 4.031 \\
bash 5.2.37   & 7.0010 & 61.29 & 2.9819 & 24.74 & 2.348 & 4.019 \\
dash (Pacman) & 2.5043 & 10.94 & 0.8096 &  7.38 & 3.093 & 1.695 \\
\bottomrule
\end{tabular}
\end{table}

\subsection{Consistency}

Standard deviations are low relative to the means in all cases.
For the Pacman \textsc{dash}, the \texttt{if} standard deviation is
\SI{10.94}{\milli\second} against a mean of \SI{2.504}{\second},
a coefficient of variation below 0.5\%.
The \texttt{case} standard deviation for the self-compiled
\textsc{dash} is the largest at \SI{71.40}{\milli\second} against
\SI{3.031}{\second}, still under 2.4\%.
The measurements are stable.
The ratios are not noise.

\subsection{The Inverse Scaling Result}

The most significant finding is not the magnitude of the overhead
but its behaviour across implementations.

Bash and dash-armv8l have nearly identical absolute \texttt{if}
times (\SI{7.001}{\second} and \SI{7.062}{\second} respectively),
despite being different interpreters.
Their \texttt{case} times are also close: \SI{2.982}{\second} and
\SI{3.031}{\second}.
The ratios are $2.348\times$ and $2.330\times$.

The Pacman \textsc{dash} is substantially faster in absolute terms:
\SI{2.504}{\second} for \texttt{if}, \SI{0.810}{\second} for
\texttt{case}.
But its ratio is $3.093\times$, higher than either of the slower
shells.

This means that a better shell implementation reduced the
\texttt{case} time more than it reduced the \texttt{if} time.
The overhead attributable to the \texttt{[} round-trip did not shrink
proportionally.
The irreducible cost of argument vector construction, built-in
dispatch, and exit-code propagation was compressed less by the
optimisation than the internal dispatch path of \texttt{case} was.

The implication is that optimising the shell does not close the
semantic gap; it exposes it more clearly.

\section{Discussion}

\subsection{What the Measurement Establishes}

These results establish three things.

First, the overhead of \texttt{if}/\texttt{[} for value-based
dispatch is not theoretical.
On the test hardware, under bash, it costs approximately
\SI{4}{\second} per 100\,000 iterations compared to \texttt{case}.
At the per-iteration level, that is roughly \SI{40}{\micro\second}
of avoidable overhead per dispatch, paid entirely to satisfy the
interface contract of a command evaluator on a problem that
does not involve commands.

Second, the overhead is consistent across implementations.
The ratio does not depend on which shell is used or how it was
compiled.
Every shell tested showed \texttt{case} faster by at least
$2.3\times$.
The structural argument predicts this: the cost is in the protocol,
not the implementation.

Third, and most importantly, optimising the shell widens the gap.
This is not a result that follows immediately from the structural
argument alone.
It requires measurement to observe.
The Pacman \textsc{dash}, the fastest shell in the test, shows
the largest ratio.
A programmer who believed that a sufficiently optimised shell would
make the distinction irrelevant is empirically wrong.

\subsection{Limitations}

Four runs per shell is a small sample.
The standard deviations are low enough that the ratios are credible,
but a larger study with process isolation, CPU pinning, and multiple
machines would produce more defensible confidence intervals.

All measurements were taken on a single architecture: 32-bit ARM
in-order pipeline.
The Cortex-A53's in-order execution makes timing more stable than
an out-of-order desktop core, but it may not represent the relative
overhead on x86\_64 or aarch64 systems.
The absolute times would differ; whether the ratios hold is an open
empirical question.

The benchmark measures the total cost of dispatch as written, not
the per-invocation cost of \texttt{[} in isolation.
Isolating the per-invocation cost would require a different
experimental design.

These limitations do not affect the primary claim: the overhead is
real, it is consistent, and it does not diminish as the
implementation improves.

\section{Conclusion}

The \texttt{if}/\texttt{[} idiom for value-based dispatch in POSIX
shell incurs a measurable and consistent overhead relative to
\texttt{case}.
On three shell configurations tested on ARM Cortex-A53 hardware,
\texttt{case} was faster by factors of $2.33\times$, $2.35\times$,
and $3.09\times$.
The overhead is structural: it follows from the execution model of
\texttt{if} and cannot be eliminated by implementation quality.
The evidence for this is in the direction of the ratio under
optimisation.
The gap does not close.
It grows.

\vspace{10pt}
\noindent\rule{\linewidth}{0.4pt}
\vspace{4pt}

{\small
\noindent
\textit{Ivan Gaydardzhiev, Independent Research, 2026.
This work is licensed under the Creative Commons Attribution 4.0
International License (CC~BY~4.0).}
}

\begin{thebibliography}{1}
\bibitem{ecnb2026}
I.~Gaydardzhiev,
\textit{Exit Codes Are Not Booleans: On the Misuse of \texttt{if}
in POSIX Shell Scripting},
Independent Research, 2026.
\end{thebibliography}

\end{document}
